%% kapa: exam=1 round=32 year=2021 period=1 form=A subject=real_estate subject_ko=부동산학원론 seq=19 number=99
%% source: Q-Net 감정평가사 기출문제 (cst00309) artlSeq=5208397 file=extracted/1008/5208397/A형/2021년도 제32회 감정평가사 1차 1교시 A형.pdf
%% conversion: kordoc 4.24.51 -> tools/gen1.py (auto); review=proofread-sonnet; flags=none
\begin{question}{99}
화폐의 시간가치계산에 관한 설명으로 옳은 것은?
\begin{choices}
  \item 연금의 현재가치계수에 일시불의 미래가치계수를 곱하면 연금의 미래가치계수가 된다.
  \item 원금균등분할상환방식에서 매 기간의 상환액을 계산할 경우 저당상수를 사용한다.
  \item 기말에 일정 누적액을 만들기 위해 매 기간마다 적립해야 할 금액을 계산할 경우 연금의 현재가치계수를 사용한다.
  \item 연금의 미래가치계수에 일시불의 현재가치계수를 곱하면 일시불의 미래가치계수가 된다.
  \item 저당상수에 연금의 현재가치계수를 곱하면 일시불의 현재가치가 된다.
\end{choices}
\end{question}
\answer{1}
